The relations to be proved can be checked pointwise on @@M0@@, so we are reduced to the case where @@M1@@ is local, with @@M2@@. The @@M3@@ are then free over @@M4@@, and the relations @@M5@@ for @@M6@@ are equivalent to the relations @@M7@@ for @@M8@@. Now we can compute the @@M9@@ using the projective resolution @@M10@@ of @@M11@@; they are the homology groups of the complex @@M12@@. Since @@M13@@ is assumed regular, this last complex is acyclic in degrees @@M14@@, whence the relations @@M15@@, for @@M16@@.

We deduce the relations @@M17@@ for @@M18@@ by induction on @@M19@@ using the exact Tor sequence associated with the exact sequence:
@@M20@@
Since for every @@M21@@ we have @@M22@@, this exact sequence and @@M23@@ give the relations @@M24@@. From the exact sequence:
@@M25@@
and @@M26@@, we deduce the relations @@M27@@. Finally, from the exact Tor sequence associated with the exact sequence:
@@M28@@
we deduce the relations @@M29@@ for @@M30@@ by induction on @@M31@@, the case @@M32@@ being evident.

b) We shall now show that the @@M33@@ are projective over @@M34@@. Since the question is local on @@M35@@, we may suppose @@M36@@ free, say @@M37@@, so that @@M38@@ is defined by a sequence @@M39@@ of elements of @@M40@@. Put @@M41@@ and let @@M42@@ be the homomorphism that sends
